\documentclass{ximera}

\input{../../../preamble.tex}


\definecolor{fvdnavy}{HTML}{071D43}
\definecolor{fvdcyan}{HTML}{20BFC6}
\definecolor{fvdcyanlight}{HTML}{EFFCFB}
\definecolor{fvdmagenta}{HTML}{F10D69}
\definecolor{fvdmagentalight}{HTML}{FFF1F7}
\definecolor{fvdtext}{HTML}{163044}





%=================================================
% Pedagogische omgevingen
% PDF: tcolorbox
% HTML: learning-box
%=================================================

\newenvironment{voorkennis}
{%
    \ifdefined\HCode
        \HCode{%
<section class="learning-box learning-box--prior">
<div class="learning-box__header">
<span class="learning-box__icon" aria-hidden="true"></span>
<span class="learning-box__title">Voorkennis</span>
</div>
<div class="learning-box__content">
        }%
        \begin{itemize}
    \else
       \begin{tcolorbox}[
    enhanced,
    breakable,
    colback=fvdcyanlight,
    colframe=fvdcyan,
    coltext=fvdtext,
    boxrule=0.5pt,
    leftrule=4pt,
    arc=4mm,
    outer arc=4mm,
    left=7mm,
    right=6mm,
    top=5mm,
    bottom=5mm,
    before skip=6mm,
    after skip=6mm,
    title=Voorkennis,
    coltitle=fvdcyan,
    fonttitle=\bfseries\Large,
    attach title to upper={\par\smallskip},
    boxed title style={
        empty,
        left=0mm,
        right=0mm,
        top=0mm,
        bottom=0mm
    },
    overlay={
        \draw[
            fvdcyan,
            line width=0.5pt
        ]
        ([xshift=42mm,yshift=-13mm]frame.north west)
        --
        ([xshift=-7mm,yshift=-13mm]frame.north east);
    }
]
        \begin{itemize}
    \fi
}
{%
    \end{itemize}
    \ifdefined\HCode
        \HCode{%
</div>
</section>
        }%
    \else
        \end{tcolorbox}
    \fi
}


\newenvironment{leerdoelen}
{%
    \ifdefined\HCode
        \HCode{%
<section class="learning-box learning-box--goals">
<div class="learning-box__header">
<span class="learning-box__icon" aria-hidden="true"></span>
<span class="learning-box__title">Leerdoelen</span>
</div>
<div class="learning-box__content">
        }%
        \begin{itemize}
    \else
\begin{tcolorbox}[
    enhanced,
    breakable,
    colback=fvdmagentalight,
    colframe=fvdmagenta,
    coltext=fvdtext,
    boxrule=0.5pt,
    leftrule=4pt,
    arc=4mm,
    outer arc=4mm,
    left=7mm,
    right=6mm,
    top=5mm,
    bottom=5mm,
    before skip=6mm,
    after skip=6mm,
    title=Leerdoelen,
    coltitle=fvdmagenta,
    fonttitle=\bfseries\Large,
    attach title to upper={\par\smallskip},
    boxed title style={
        empty,
        left=0mm,
        right=0mm,
        top=0mm,
        bottom=0mm
    },
    overlay={
        \draw[
            fvdmagenta,
            line width=0.5pt
        ]
        ([xshift=42mm,yshift=-13mm]frame.north west)
        --
        ([xshift=-7mm,yshift=-13mm]frame.north east);
    }
]
        \begin{itemize}
    \fi
}
{%
    \end{itemize}
    \ifdefined\HCode
        \HCode{%
</div>
</section>
        }%
    \else
        \end{tcolorbox}
    \fi
}


\addPrintStyle{../../..}

\title{Oefeningen — Symmetrie}
\author{Ingmar Herreman}

\begin{document}

% FVD_CHAPTER_DATA_START
% Automatisch gegenereerd uit index.tex.
% Niet handmatig aanpassen.
\fvdchapterdata{5}{Het gedrag van functies beschrijven}{4}
% FVD_CHAPTER_DATA_END




\fvdchapterdata{5}{Het gedrag van functies beschrijven}{4}

\maketitle
\FVDbeginExercises

\begin{opmerking}
De oefeningen met \(\blacklozenge\) zijn essentieel.
Je moet symmetrie zowel grafisch als algebraïsch kunnen analyseren en beide
voorstellingen met elkaar verbinden.
\end{opmerking}

\FVDsection{Basis — even en oneven herkennen}

\begin{oefening}[Functiewaarden en symmetrie]{\oefB\ESS}

\begin{question}
Een even functie voldoet aan \(f(4)=7\). Bepaal \(f(-4)\).
\begin{oplossing}
Omdat \(f\) even is:
\[
f(-4)=f(4)=7.
\]
\end{oplossing}
\end{question}

\begin{question}
Een oneven functie voldoet aan \(g(3)=-5\). Bepaal \(g(-3)\).
\begin{oplossing}
Omdat \(g\) oneven is:
\[
g(-3)=-g(3)=5.
\]
\end{oplossing}
\end{question}

\begin{question}
Een oneven functie heeft \(0\) in haar domein. Bepaal \(g(0)\).
\begin{oplossing}
\[
g(0)=-g(0)\Rightarrow g(0)=0.
\]
\end{oplossing}
\end{question}

\end{oefening}

\begin{oefening}[Algebraïsch onderzoeken]{\oefB\ESS}

Onderzoek of de functie even, oneven of geen van beide is.

\begin{question}
\[
f(x)=x^4-3x^2+1.
\]
\begin{oplossing}
\[
f(-x)=x^4-3x^2+1=f(x).
\]
Dus \(f\) is even.
\end{oplossing}
\end{question}

\begin{question}
\[
g(x)=x^5-2x^3+x.
\]
\begin{oplossing}
\[
g(-x)=-x^5+2x^3-x=-g(x).
\]
Dus \(g\) is oneven.
\end{oplossing}
\end{question}

\begin{question}
\[
h(x)=x^2+2x+1.
\]
\begin{oplossing}
\[
h(-x)=x^2-2x+1.
\]
Dit is noch \(h(x)\), noch \(-h(x)\). Dus \(h\) is geen van beide.
\end{oplossing}
\end{question}

\end{oefening}

\begin{oefening}[Punten reconstrueren]{\oefB\ESS}

\begin{question}
De grafiek van een even functie bevat
\[
(1,3),\quad(2,-1),\quad(4,5).
\]
Geef drie andere punten die zeker op de grafiek liggen.
\begin{oplossing}
\[
(-1,3),\quad(-2,-1),\quad(-4,5).
\]
\end{oplossing}
\end{question}

\begin{question}
De grafiek van een oneven functie bevat
\[
(1,3),\quad(2,-1),\quad(4,5).
\]
Geef drie andere punten die zeker op de grafiek liggen.
\begin{oplossing}
\[
(-1,-3),\quad(-2,1),\quad(-4,-5).
\]
\end{oplossing}
\end{question}

\end{oefening}

\begin{oefening}[Rationale functies]{\oefB\ESS}

Onderzoek de symmetrie.

\begin{question}
\[
f(x)=\frac{x}{x^2+1}.
\]
\begin{oplossing}
Het domein is \(\mathbb R\).
\[
f(-x)=\frac{-x}{x^2+1}=-f(x).
\]
Dus \(f\) is oneven.
\end{oplossing}
\end{question}

\begin{question}
\[
g(x)=\frac{x^2}{x^2+1}.
\]
\begin{oplossing}
\[
g(-x)=\frac{x^2}{x^2+1}=g(x).
\]
Dus \(g\) is even.
\end{oplossing}
\end{question}

\end{oefening}

\FVDsection{Verdieping — symmetrie analyseren}

\begin{oefening}[Het domein telt mee]{\oefU\ESS}

Beschouw
\[
f:[0,+\infty[\to\mathbb R:x\mapsto x^2.
\]

Een leerling zegt:
``\(f\) is even, want \(f(-x)=f(x)\).''

Analyseer de uitspraak.

\begin{oplossing}
De uitspraak is fout voor deze functie.
Het domein is \([0,+\infty[\) en is niet symmetrisch ten opzichte van \(0\).
Als bijvoorbeeld \(2\) in het domein ligt, ligt \(-2\) er niet in.
De algebraïsche vorm \(x^2\) alleen volstaat dus niet; het domein hoort bij de functie.
\end{oplossing}

\end{oefening}

\begin{oefening}[Van nulwaarde naar een nieuwe nulwaarde]{\oefU\ESS}

\begin{question}
Een even functie heeft \(x=5\) als nulwaarde. Welke andere nulwaarde volgt
noodzakelijk?
\begin{oplossing}
\[
f(5)=0\Rightarrow f(-5)=f(5)=0.
\]
Dus \(-5\) is eveneens een nulwaarde.
\end{oplossing}
\end{question}

\begin{question}
Een oneven functie heeft \(x=5\) als nulwaarde. Welke andere nulwaarde volgt?
\begin{oplossing}
\[
f(-5)=-f(5)=0.
\]
Dus \(-5\) is eveneens een nulwaarde.
\end{oplossing}
\end{question}

\end{oefening}

\begin{oefening}[Grafische informatie koppelen]{\oefU\ESS}

Een even functie is voor \(x>0\) stijgend.

\begin{question}
Wat kun je besluiten over het verloop voor \(x<0\)?
\begin{oplossing}
Door de spiegeling in de \(y\)-as is de linkerhelft gespiegeld.
Als de functie rechts stijgt wanneer we van links naar rechts bewegen,
dan daalt ze links wanneer we van links naar rechts bewegen.
\end{oplossing}
\end{question}

\begin{question}
Leg uit waarom deze conclusie uit de symmetrie volgt.
\begin{oplossing}
Bij een even functie hebben \(x\) en \(-x\) dezelfde functiewaarde.
De volgorde van de overeenkomstige punten wordt links van de \(y\)-as omgekeerd,
waardoor stijgen rechts overeenkomt met dalen links.
\end{oplossing}
\end{question}

\end{oefening}

\begin{oefening}[Is drie punten controleren voldoende?]{\oefU\ESS}

Een leerling onderzoekt een functie en vindt:
\[
f(-1)=f(1),\qquad f(-2)=f(2),\qquad f(-3)=f(3).
\]
Hij besluit dat \(f\) even is.

Is dit voldoende? Leg uit.

\begin{oplossing}
Nee. Voor evenheid moet
\[
f(-x)=f(x)
\]
gelden voor elke \(x\) uit het domein. Drie controles ondersteunen een vermoeden,
maar bewijzen de algemene eigenschap niet.
\end{oplossing}

\end{oefening}

\begin{oefening}[Voorschrift en grafiek verbinden]{\oefU\ESS}

Gegeven:
\[
f(x)=\frac{x^3}{x^2-4}.
\]

\begin{question}
Bepaal het domein.
\begin{oplossing}
\[
\operatorname{dom}f=\mathbb R\setminus\{-2,2\}.
\]
\end{oplossing}
\end{question}

\begin{question}
Toon algebraïsch aan dat \(f\) oneven is.
\begin{oplossing}
\[
f(-x)=\frac{-x^3}{x^2-4}=-f(x).
\]
Het domein is bovendien symmetrisch ten opzichte van \(0\).
\end{oplossing}
\end{question}

\begin{question}
Als \((3,27/5)\) op de grafiek ligt, welk punt volgt onmiddellijk uit de symmetrie?
\begin{oplossing}
\[
(-3,-27/5).
\]
\end{oplossing}
\end{question}

\end{oefening}

\FVDsection{Gevorderd — redeneren met symmetrie}

\begin{oefening}[Kan een functie beide zijn?]{\oefG}

Kan een functie tegelijk even en oneven zijn?
Onderzoek de voorwaarden
\[
f(-x)=f(x)
\qquad\text{en}\qquad
f(-x)=-f(x).
\]

\begin{oplossing}
Als beide gelden, dan
\[
f(x)=-f(x)
\]
voor elke \(x\) uit het domein. Dus
\[
2f(x)=0
\quad\Rightarrow\quad
f(x)=0.
\]
Op een domein dat symmetrisch is ten opzichte van \(0\) is de nulfunctie dus
zowel even als oneven.
\end{oplossing}

\end{oefening}

\begin{oefening}[Ontwerp een functie]{\oefG}

Geef:
\begin{enumerate}
  \item een even veeltermfunctie met nulwaarden \(-2\) en \(2\);
  \item een oneven veeltermfunctie met nulwaarden \(-2\), \(0\) en \(2\);
  \item een functie met dezelfde drie nulwaarden die noch even, noch oneven is.
\end{enumerate}

Motiveer telkens.

\begin{oplossing}
Mogelijke antwoorden zijn:
\[
f(x)=x^2-4,
\]
die even is;

\[
g(x)=x(x^2-4)=x^3-4x,
\]
die oneven is;

en bijvoorbeeld
\[
h(x)=x(x-2)(x+2)(x-1),
\]
die dezelfde drie gevraagde nulwaarden bevat maar in het algemeen noch even,
noch oneven is.

Andere correcte voorbeelden zijn mogelijk.
\end{oplossing}

\end{oefening}

\begin{oefening}[GeoGebra-onderzoek]{\oefG}

Onderzoek in GeoGebra de familie
\[
f_a(x)=x^4+a x^2
\]
en
\[
g_a(x)=x^3+a x.
\]

Gebruik een schuifknop voor \(a\).

\begin{enumerate}
  \item Welke symmetrie blijft bij elke waarde van \(a\) behouden?
  \item Toon dit ook algebraïsch aan.
  \item Onderzoek hoe nulwaarden en extrema kunnen veranderen terwijl de symmetrie behouden blijft.
  \item Leg uit waarom symmetrie een kenmerk is dat onafhankelijk kan blijven van andere
        grafische kenmerken.
\end{enumerate}

\begin{oplossing}
Voor elke \(a\) bevat \(f_a\) alleen even machten:
\[
f_a(-x)=f_a(x),
\]
dus \(f_a\) blijft even.

Voor \(g_a\):
\[
g_a(-x)=-x^3-a x=-g_a(x),
\]
dus \(g_a\) blijft oneven.

Nulwaarden en extrema kunnen met \(a\) veranderen, terwijl de algebraïsche
symmetrievoorwaarden behouden blijven.
\end{oplossing}

\end{oefening}



\end{document}
